\subsection{附加于有限生成模的除子类}

我们保持第 2 至第 6 节的记号与假设。回忆：C(A)（或简作 C）表示 A 的除子类群，即 D(A) 关于主除子所成子群的商。对每个除子 \(d \in D\) ，我们用 \(c(d)\) 表示它在 C 中的类。

命题 15. 设 M 是有限生成 A-模。存在 M 的自由子模 L，使 M/L 是扭模，且 C 的元素 \(c(\chi(M/L))\) 与所选的自由子模 L 无关。

我们写 \(S = A - \{0\}\) 并令 \(V = S^{-1}M = M \otimes_{A} K\) ；若 \(n\) 是 \(V\) 在 \(K\) 上的秩，则存在 \(n\) 个元素 \(e_i (1 \leqslant i \leqslant n)\) 属于 \(M\) ，它们的典范像在 \(V\) 中构成 \(V\) 的基；这些元素在 \(M\) 中显然线性无关，从而生成自由子模 \(L\) ，它含于 \(M\) ，满足 \(S^{-1}(M/L) = S^{-1}M/S^{-1}L = 0\) ，因此 \(M/L\) 是扭模。

现在设 \(L_{1}\) 是 M 的另一个秩为 n 的自由子模。~由于 \(S^{-1}L = S^{-1}L_{1}\) ，存在 \(s \in S\) 使 \(sL_{1} \subset L\) ；因此我们只需限于证明：若 \(L_{1} \subset L_{2}\) 是 M 的两个秩为 n 的自由子模，则

\[
c (\chi (\mathbf {M} / \mathbf {L} _ {1})) = c (\chi (\mathbf {M} / \mathbf {L} _ {2})).
\]

现在，\(\chi (\mathbf{M} / \mathbf{L}_1) = \chi (\mathbf{M} / \mathbf{L}_2) + \chi (\mathbf{L}_2 / \mathbf{L}_1)\) ，且由第 6 节命题 14 的推论可知 \(\chi (\mathbf{L}_2 / \mathbf{L}_1)\) 是主除子，因此

\[
c (\chi (\mathbf {L} _ {2} / \mathbf {L} _ {1})) = 0.
\]

元素 \(c(\chi(\mathbf{M} / \mathbf{L}))\) 以下记作 \(-c(\mathbf{M})\) ；我们称 \(c(\mathbf{M})\) 为附加于 \(\mathbf{M}\) 的除子类。

命题 16. (i) 设 \(0 \to \mathbf{M}_1 \xrightarrow{f} \mathbf{M}_2 \xrightarrow{g} \mathbf{M}_3 \to 0\) 是有限生成 A-模的正合列。则

\[
c (\mathbf {M _ {2}}) = c (\mathbf {M _ {1}}) + c (\mathbf {M _ {3}}).
\]

\begin{enumerate}
\def\labelenumi{(\roman{enumi})}
\setcounter{enumi}{1}
\item
  若存在从 \(\mathbf{M}_1\) 到 \(\mathbf{M}_2\) 的伪同构，则 \(c(\mathbf{M}_1) = c(\mathbf{M}_2)\) 。
\item
  若 \(\mathbf{T}\) 是扭模，则 \(c(\mathbf{T}) = -c(\chi(\mathbf{T}))\) 。
\item
  若 \(\mathfrak{a} \neq 0\) 是 \(\mathbf{K}\) 的分式理想，则
\end{enumerate}

\[
c (\mathfrak {a}) = c (\operatorname{div} (\mathfrak {a})).
\]

\begin{enumerate}
\def\labelenumi{(\alph{enumi})}
\setcounter{enumi}{21}
\tightlist
\item
  若 \(\mathbf{L}\) 是自由 A-模，则 \(c(\mathbf{L}) = 0\) 。
\end{enumerate}

为证 (i)，考虑自由子模 \(\mathbf{L}_1\) （相应地，\(\mathbf{L}_3\) ），它含于 \(\mathbf{M}_1\) （相应地，\(\mathbf{M}_3\) ），使 \(\mathbf{M}_1 / \mathbf{L}_1\) （相应地，\(\mathbf{M}_3 / \mathbf{L}_3\) ）是扭模。由于 \(\mathbf{L}_3\) 自由且 \(g\) 是满射，在 \(\overline{g}^{-1} (\mathbf{L}_3)\) 中存在一个自由补 \(\mathbf{L}_{23}\) （ \(\operatorname {Ker}(g)\) 的），它与 \(\mathbf{L}_3\) 同构 （《代数》第 II 章，§ 1，第 11 节，命题 21）；但 \(\operatorname {Ker}(g) = f(\mathbf{M}_1)\)

它包含 \(f(\mathbf{L}_1) = \mathbf{L}_{12}\) ，后者因 \(f\) 是单射而是自由的。\(\mathbf{L}_2 = \mathbf{L}_{12} + \mathbf{L}_{23}\) 这个和是直和，因此 \(\mathbf{L}_2\) 是 \(\mathbf{M}_2\) 的自由子模。此外还有交换图：

\[
\begin{array}{c} 0 \longrightarrow \mathrm{L} _ {1} \longrightarrow \mathrm{L} _ {2} \longrightarrow \mathrm{L} _ {3} \longrightarrow 0 \\ \Biggl \downarrow \qquad \qquad \qquad \Biggl \downarrow \qquad \qquad \Biggl \downarrow \\ 0 \longrightarrow \mathrm{M} _ {1} \xrightarrow [ f ]{} \mathrm{M} _ {2} \xrightarrow [ g ]{} \mathrm{M} _ {3} \longrightarrow 0 \end{array}
\]

其中各行正合，各竖直箭均为单射。因此我们由蛇形图（第 I 章，§ 1，第 4 节，命题 2）得到正合列：

\[
0 \rightarrow \mathrm{M} _ {1} / \mathrm{L} _ {1} \rightarrow \mathrm{M} _ {2} / \mathrm{L} _ {2} \rightarrow \mathrm{M} _ {3} / \mathrm{L} _ {3} \rightarrow 0.
\]

由于 \(\mathbf{M}_1 / \mathbf{L}_1\) 与 \(\mathbf{M}_3 / \mathbf{L}_3\) 都是扭模，这一正合列首先表明 \(\mathbf{M}_2 / \mathbf{L}_2\) 亦然，继而由第 5 节命题 10 得：

\[
\chi (\mathrm{M} _ {2} / \mathrm{L} _ {2}) = \chi (\mathrm{M} _ {1} / \mathrm{L} _ {1}) + \chi (\mathrm{M} _ {3} / \mathrm{L} _ {3})
\]

这就证明了 (i)。

断言 (iii) 与 (v) 由定义立得。我们证明 (ii)。于是设 \(f \colon \mathbf{M}_1 \to \mathbf{M}_2\) 是伪同构，\(\mathbf{L}_1\) 是 \(\mathbf{M}_1\) 的自由子模，使 \(\mathbf{M}_1 / \mathbf{L}_1\) 是扭模。我们置 \(\mathbf{L}_2 = f(\mathbf{L}_1)\) ；由于 \(\operatorname{Ker}(f)\) 是伪零的，它是扭模，故 \(\operatorname{Ker}(f) \cap \mathbf{L}_1 = 0\) ，从而 \(\mathbf{L}_2\) 是自由的。设 \(\bar{f} \colon \mathbf{M}_1 / \mathbf{L}_1 \to \mathbf{M}_2 / \mathbf{L}_2\) 是由 \(f\) 取商导出的同态；\(\operatorname{Ker}(\bar{f})\) 同构于 \(\operatorname{Ker}(f)\) ，\(\operatorname{Coker}(\bar{f})\) 同构于 \(\operatorname{Coker}(f)\) ，故 \(\bar{f}\) 是伪同构；此外

\[
\operatorname{Coker} (\bar {f}) = \mathrm{M} _ {2} / f (\mathrm{M} _ {1})
\]

是扭模，\(f(\mathbf{M}_1) / \mathbf{L}_2 = \bar{f} (\mathbf{M}_1 / \mathbf{L}_1)\) 亦然，故 \(\mathbf{M}_2 / \mathbf{L}_2\) 是扭模，于是由第 5 节命题 10 (ii) 得

\[
\chi (\mathrm{M} _ {1} / \mathrm{L} _ {1}) = \chi (\mathrm{M} _ {2} / \mathrm{L} _ {2}).
\]

最后尚须证明 (iv)。设 \(x \in \mathbf{K}^*\) 满足 \(a \subset xA\) 。考虑正合列 \(0 \to a \to xA \to xA / a \to 0\) ，我们得到

\[
c (\mathfrak {a}) = c (x \mathrm{A}) - c (x \mathrm{A} / \mathfrak {a}) = - c (x \mathrm{A} / \mathfrak {a})
\]

这是依据 (i) 与 (v)。但 \(x\mathrm{A} / \mathfrak{a}\) 同构于 \(\mathrm{A} / x^{-1}\mathfrak{a}\) ，于是依据 (iii)，

\[
c (x \mathrm{A} / \mathrm{a}) = - c (\chi (\mathrm{A} / x ^ {- 1} \mathrm{a})) = - c (\operatorname{div} (x ^ {- 1} \mathrm{a})) = - c (\operatorname{div} (\mathrm{a}))
\]

（第 5 节，命题 12）。证毕。

当 M 是 V 关于 A 的格时，\(\chi(\mathrm{M}/\mathrm{L}) = -\chi(\mathrm{M}, \mathrm{L})\) （第 6 节，命题 14）；设 \((e_{i})_{1 \leqslant i \leqslant n}\) 是 L 的基，\(e = e_{1} \wedge e_{2} \wedge \cdots \wedge e_{n}\) ，且

\(M_{w} = a.e\) （第 6 节的记号）；则 \(\chi(M, L) = \text{div}(a)\) ，从而 \(c(M) = c(\text{div}(a))\) ，这推广了命题 16 (v)。

推论 1. 设 \(0 \to \mathbf{M}_n \xrightarrow{u} \mathbf{M}_{n-1} \to \cdots \to \mathbf{M}_0 \to 0\) 是有限生成 A-模的正合列。则

\[
\sum_ {i = 0} ^ {n} (- 1) ^ {i} c (\mathbf {M} _ {i}) = 0.
\]

我们对 \(n\) 作归纳，\(n = 2\) 的情形即命题 16 (i)。若 \(\mathbf{M}_{n-1}' = \operatorname{Coker}(u)\) ，则有两个正合列：

\[
0 \rightarrow \mathrm{M} _ {n} \rightarrow \mathrm{M} _ {n - 1} \rightarrow \mathrm{M} _ {n - 1} ^ {\prime} \rightarrow 0
\]

\[
0 \rightarrow \mathrm{M} _ {n - 1} ^ {\prime} \rightarrow \mathrm{M} _ {n - 2} \rightarrow \dots \rightarrow \mathrm{M} _ {0} \rightarrow 0.
\]

第一个正合列表明 \(\mathbf{M}_{n - 1}^{\prime}\) 是有限生成的，且归纳假设给出

\[
(- 1) ^ {n - 1} c (\mathbf {M} _ {n - 1} ^ {\prime}) + \sum_ {i = 0} ^ {n - 2} (- 1) ^ {i} c (\mathbf {M} _ {i}) = 0
\]

且

\[
c (\mathbf {M} _ {n - 1} ^ {\prime}) = c (\mathbf {M} _ {n - 1}) - c (\mathbf {M} _ {n}),
\]

由此得推论。

正合列

\[
0 \rightarrow \mathrm{L} _ {n} \rightarrow \mathrm{L} _ {n - 1} \rightarrow \dots \rightarrow \mathrm{L} _ {0} \rightarrow \mathrm{E} \rightarrow 0
\]

其中诸 \(L_{i}\) （ \(0 \leqslant i \leqslant n\) ）是有限生成的自由 A-模，这样的正合列称为 A-模 E 的有限自由分解。

推论 2. 若 A 的除性分式理想 \(\mathfrak{a} \neq 0\) 容许一个有限自由分解，则它是主理想。

事实上，我们对 \(\mathfrak{a}\) 的一个有限自由分解应用推论 1：

\[
0 \rightarrow \mathrm{L} _ {n} \rightarrow \mathrm{L} _ {n - 1} \rightarrow \dots \rightarrow \mathrm{L} _ {0} \rightarrow \mathfrak {a} \rightarrow 0.
\]

由命题 16 (v)，\(c(a) = 0\) ，从而由命题 16 (iv)，\(\text{div}(a)\) 是主理想；由于假设 \(a\) 是除性的，故它是主理想（§ 1，第 1 节）。

推论 3. 若 A 的每个 \(\neq 0\) 除性理想都容许有限自由分解，则 A 是因子分解整环。

这是推论 2 与 § 3 第 1 节定义 1 的直接推论。

* 我们以后将看到，正则局部环满足推论

3 的假设，从而是因子分解整环。*

若 M 是有限生成 A-模，我们将用 \(r(\mathbf{M})\) 表示它的秩（回忆它是 \(\mathbf{M}_{(\mathbf{K})} = \mathbf{M} \otimes_{\mathbf{A}} \mathbf{K}\) 在 K 上的秩）；若 \(0 \to M_{1} \to M_{2} \to M_{3} \to 0\) 是有限生成 A-模的正合列，则序列

\[
0 \rightarrow (\mathrm{M} _ {1}) _ {(\mathbf {K})} \rightarrow (\mathrm{M} _ {2}) _ {(\mathbf {K})} \rightarrow (\mathrm{M} _ {3}) _ {(\mathbf {K})} \rightarrow 0
\]

也正合，因此 \(r(\mathbf{M}_2) = r(\mathbf{M}_1) + r(\mathbf{M}_3)\) 。我们写

\[
\gamma (\mathbf {M}) = (r (\mathbf {M}), c (\mathbf {M})) \in \mathbf {Z} \times \mathbf {C} (\mathbf {A});
\]

于是 γ 满足命题 16 的性质 (i)，且当 M 是伪零模时 γ(M) = 0（因为 M 是扭模）。存在从 F(A) 到 Z × C(A) 的唯一映射，仍记作 γ，使得对每个有限生成 A-模 M 有 γ(M) = γ(cl(M))。我们将看到，上述性质本质上刻画了 γ：

命题 17. 设 G 是以加法写出的交换群，\(\phi\) 是从有限生成 A-模的类组成的集 \(F(A)\) 到 G 的映射；对每个有限生成 A-模 M，我们也滥用语言写 \(\phi(M) = \phi(cl(M))\) 。设下列条件满足：

\begin{enumerate}
\def\labelenumi{(\arabic{enumi})}
\item
  若 \(0 \to \mathbf{M}_1 \to \mathbf{M}_2 \to \mathbf{M}_3 \to 0\) 是有限生成 A-模的正合列，则 \(\phi(\mathbf{M}_2) = \phi(\mathbf{M}_1) + \phi(\mathbf{M}_3)\) 。
\item
  若 \(\mathbf{T}\) 是伪零的，则 \(\phi (\mathbf{T}) = 0\) 。
\end{enumerate}

则存在唯一同态 \(\theta: \mathbf{Z} \times \mathbf{C} \to \mathbf{G}\) 使 \(\phi = \theta \circ \gamma\) 。

由命题 16 (iv)，\(\mathbf{Z} \times \mathbf{C}\) 的每个元素对某个适当的有限生成 A-模 M 都形如 \((r(\mathbf{M}), c(\mathbf{M}))\) ；由此得 \(\theta\) 的唯一性。我们对 \(-\phi\) 在 \(T(\mathbf{A})\) 上的限制应用第 5 节的命题 11 ：于是存在同态 \(\theta_0: \mathbf{D} \to \mathbf{G}\) 使

\[
- \phi (\mathbf {T}) = \theta_ {0} (\chi (\mathbf {T}))
\]

对每个有限生成扭 A-模 T 成立。设 \(x\) 是 A 的非零元素；把性质 (1) 应用于正合列：

\[
0 \longrightarrow \mathrm{A} \xrightarrow {h _ {x}} \mathrm{A} \longrightarrow \mathrm{A} / x \mathrm{A} \longrightarrow 0
\]

其中 \(h_{x}\) 是乘以 x 的映射，我们得到 \(\phi(\mathbf{A}/x\mathbf{A}) = 0\) ，从而 \(\theta_{0}(\mathrm{div}(x)) = 0\) 。取商后，\(\theta_{0}\) 便定义了一个同态 \(\theta_{1}: C \to G\) ，且对每个扭 A-模 T 有 \(\phi(T) = \theta_{1}(c(T))\) 。我们现在证明同态 \(\theta\) （由 \(\theta(n, z) = n.\phi(\mathbf{A}) + \theta_{1}(z)\) 定义）解决了问题。为此，对每个有限生成 A-模 M 写 \(\phi'(\mathbf{M}) = \phi(\mathbf{M}) - \theta(\gamma(\mathbf{M}))\) ；显然当把 \(\phi\) 换成 \(\phi'\) 时条件 (1) 仍满足。此外，当 M 是扭模或自由模时 \(\phi'(\mathbf{M}) = 0\) ；但对每个有限生成 A-模 M，存在 M 的自由子模 L 使 M/L 是扭模（命题 15），故性质 (1) 表明对每个有限生成 A-模 M 都有 \(\phi'(\mathbf{M}) = 0\) 。

* 用 Abel 范畴的语言，命题 17 表明 \(\mathbf{Z} \times \mathrm{C}(\mathbf{A})\) 典范同构于商范畴 \(\mathcal{F}/\mathcal{F}'\) 的 Grothendieck 群，其中 \(\mathcal{F}\) 是有限生成 A-模的范畴，\(\mathcal{F}'\) 是 \(\mathcal{F}\) 中由伪零模组成的全子范畴。*

